\chapter{Requisiti e modello delle minacce}
\label{cap:requisiti}

\section{Attori e casi d'uso}
\todo{Sviluppatore, verificatore, ente certificatore, consumatore del software.}

\section{Requisiti funzionali e non funzionali}
\begin{table}[htbp]
  \centering
  \caption{Elenco dei requisiti.}
  \label{tab:requisiti}
  \begin{tabularx}{\linewidth}{@{}llX@{}}
    \toprule
    \textbf{ID} & \textbf{Tipo} & \textbf{Descrizione}\\
    \midrule
    RF-01 & Funzionale & \todo{Generazione della \sbom{} a partire dal progetto.}\\
    RF-02 & Funzionale & \todo{Emissione e verifica del certificato di sicurezza.}\\
    RNF-01 & Sicurezza & \todo{Integrità e non ripudio dei certificati.}\\
    \bottomrule
  \end{tabularx}
\end{table}

\section{Modello delle minacce}
\todo{Metodologia adottata (ad esempio STRIDE), asset, avversari, superfici di attacco e contromisure.}
